%% Pages liminaires (provisoires) : titre, légende des mains et des signes.
\thispagestyle{empty}
\vspace*{60mm}
\begin{center}
{\fontsize{30}{36}\selectfont\addfontfeature{LetterSpace=8}ARISTOPHANE}\par\vspace{9mm}
{\fontsize{20}{26}\selectfont\itshape Les Thesmophories}\par\vspace{4mm}
{\capsgr{\fontsize{14}{18}\selectfont ΘΕΣΜΟΦΟΡΙΑΖΟΥΣΑΙ}}\par\vspace{18mm}
{\large Édition du \textit{Ravennas} (Ravenne, Biblioteca Classense, ms.~429),\\[1mm]
texte établi, traduction et commentaire}\par
\end{center}
\clearpage
\entetes{}{}
\section*{Légende}
\subsection*{Les mains du \textit{Ravennas}}
\begin{tabular}{@{}ll@{}}
\textcolor{mainR1}{\rule{8mm}{2mm}} & R\textsuperscript{1} : copiste du texte\\
\textcolor{mainR1c}{\rule{8mm}{2mm}} & R\textsuperscript{1c} : corrections du copiste lui-même\\
\textcolor{mainR2}{\rule{8mm}{2mm}} & R\textsuperscript{2} : correcteur ancien\\
\textcolor{mainRsch}{\rule{8mm}{2mm}} & main des scholies (si distincte)\\
\textcolor{mainRrec}{\rule{8mm}{2mm}} & main récente\\
\end{tabular}
\par\medskip
\textit{(Légende provisoire : elle sera arrêtée au terme de l'étude graphologique.)}
\subsection*{Signes de la transcription diplomatique}
\begin{tabular}{@{}ll@{}}
⟦α⟧ & lettres grattées \\
\Dras{αβ} & lettres écrites sur grattage \\
α̣ & lecture douteuse \\
{[\,··\,]} & lettres illisibles \\
κ(αί) & abréviation développée \\
\end{tabular}
\clearpage
